
\documentclass{article}
\usepackage{colortbl}
\usepackage{makecell}
\usepackage{multirow}
\usepackage{supertabular}

\begin{document}

\newcounter{utterance}

\centering \large Interaction Transcript for game `cladder', experiment `full\_v1.5\_default', episode 1290 with refineonly{-}Qwen3.5{-}9B.
\vspace{24pt}

{ \footnotesize  \setcounter{utterance}{1}
\setlength{\tabcolsep}{0pt}
\begin{supertabular}{c@{$\;$}|p{.15\linewidth}@{}p{.15\linewidth}p{.15\linewidth}p{.15\linewidth}p{.15\linewidth}p{.15\linewidth}}
    \# & \multicolumn{2}{c}{Player} && \multicolumn{2}{c}{Game Master} \\
    \hline

    \theutterance \stepcounter{utterance}  
    & & & \multicolumn{4}{p{0.6\linewidth}}{
        \cellcolor[rgb]{0.9,0.9,0.9}{
            \makecell[{{p{\linewidth}}}]{
                \texttt{\tiny{[P1$\langle$GM]}}
                \texttt{You are an expert in causal inference. The following question is not a typical commonsense query, but rather a meticulously designed question created by a professor specializing in causal inference, intended to assess the students' mastery of the course content.} \\
\texttt{Imagine a self{-}contained, hypothetical world with only the following conditions, and without any unmentioned factors or causal relationships: Gender has a direct effect on smoking and lung cancer. Smoking has a direct effect on tar deposit. Tar deposit has a direct effect on lung cancer. Gender is unobserved. For nonsmokers, the probability of high tar deposit is 36\%. For smokers, the probability of high tar deposit is 84\%. For nonsmokers and with no tar deposit, the probability of lung cancer is 12\%. For nonsmokers and with high tar deposit, the probability of lung cancer is 44\%. For smokers and with no tar deposit, the probability of lung cancer is 15\%. For smokers and with high tar deposit, the probability of lung cancer is 50\%. The overall probability of smoking is 67\%. Will smoking increase the chance of lung cancer?} \\
\texttt{Start your answer with 'Yes' or 'No', followed by additional reasoning or evidence to support your explanation.} \\
            }
        }
    }
     \\ \\

    \theutterance \stepcounter{utterance}  
    & \multicolumn{4}{p{0.6\linewidth}}{
        \cellcolor[rgb]{0.561,0.737,0.561}{
            \makecell[{{p{\linewidth}}}]{
                \texttt{\tiny{[P1$\rangle$GM]}}
                \texttt{Yes} \\
\\ 
\texttt{**Reasoning:**} \\
\\ 
\texttt{To determine if smoking increases the chance of lung cancer, we must calculate the Average Causal Effect (ACE), defined as \$P(L|Do(S=1)) {-} P(L|Do(S=0))\$. In this world, the causal structure is \$G \textbackslash{}to S\$, \$G \textbackslash{}to L\$, \$S \textbackslash{}to T\$, and \$T \textbackslash{}to L\$. Since we are intervening on Smoking (\$Do(S)\$), we block the backdoor path through Gender (\$G\$) and only consider the effect of Smoking on Tar Deposit, which subsequently affects Lung Cancer.} \\
\\ 
\texttt{First, we calculate the probability of Lung Cancer for the population forced to be smokers (\$Do(S=1)\$):} \\
\texttt{\$\$P(L|Do(S=1)) = P(L|S=1, T=High)P(T=High|S=1) + P(L|S=1, T=Low)P(T=Low|S=1)\$\$} \\
\texttt{\$\$P(L|Do(S=1)) = (0.50 \textbackslash{}times 0.84) + (0.15 \textbackslash{}times 0.16) = 0.42 + 0.024 = 0.444\$\$} \\
\\ 
\texttt{Next, we calculate the probability of Lung Cancer for the population forced to be nonsmokers (\$Do(S=0)\$):} \\
\texttt{\$\$P(L|Do(S=0)) = P(L|S=0, T=High)P(T=High|S=0) + P(L|S=0, T=Low)P(T=Low|S=0)\$\$} \\
\texttt{\$\$P(L|Do(S=0)) = (0.44 \textbackslash{}times 0.36) + (0.12 \textbackslash{}times 0.64) = 0.1584 + 0.0768 = 0.2352\$\$} \\
\\ 
\texttt{Comparing the two probabilities:} \\
\texttt{\$\$ACE = 0.444 {-} 0.2352 = 0.2088\$\$} \\
\\ 
\texttt{Since the probability of lung cancer is higher under the intervention of smoking (\$0.444\$) than under the intervention of nonsmoking (\$0.2352\$), smoking increases the chance of lung cancer.} \\
            }
        }
    }
    & & \\ \\

    \theutterance \stepcounter{utterance}  
    & & & \multicolumn{2}{p{0.3\linewidth}}{
        \cellcolor[rgb]{0.9,0.9,0.9}{
            \makecell[{{p{\linewidth}}}]{
                \texttt{\tiny{[GM$|$GM]}}
                \texttt{Yes} \\
            }
        }
    }
    & & \\ \\

    \theutterance \stepcounter{utterance}  
    & & & \multicolumn{2}{p{0.3\linewidth}}{
        \cellcolor[rgb]{0.9,0.9,0.9}{
            \makecell[{{p{\linewidth}}}]{
                \texttt{\tiny{[GM$|$GM]}}
                \texttt{yes} \\
            }
        }
    }
    & & \\ \\

    \theutterance \stepcounter{utterance}  
    & & & \multicolumn{2}{p{0.3\linewidth}}{
        \cellcolor[rgb]{0.9,0.9,0.9}{
            \makecell[{{p{\linewidth}}}]{
                \texttt{\tiny{[GM$|$GM]}}
                \texttt{game\_result = WIN} \\
            }
        }
    }
    & & \\ \\

\end{supertabular}
}

\end{document}
